\section{Review Methodology}
\label{sec:methodology}

This Systematic Literature Review is conducted in strict adherence to evidence-based software engineering guidelines and informed by PRISMA reporting principles. The methodological protocol encompasses search strategy definition, study selection via explicit inclusion and exclusion criteria, quality appraisal, and structured data extraction.

\subsection{Search strategy and electronic databases}
We queried four leading electronic bibliographic databases indexing premier software engineering literature: IEEE Xplore Digital Library, ACM Digital Library, ScienceDirect (Elsevier), and SpringerLink. To capture seminal and emerging literature that might be indexed outside these primary databases, supplementary backward and forward snowballing was conducted on key surveys and conference proceedings (e.g., IEEE RE, REFSQ, WER).

The search string was structured across three conceptual facets connected with Boolean operators: (1) \textit{Facet 1 (Intervention)}, covering terms such as \texttt{"prototyping"}, \texttt{"prototype"}, \texttt{"mock-up"}, \texttt{"wireframe"}, \texttt{"paper prototype"}, and \texttt{"interactive prototype"}; (2) \textit{Facet 2 (RE Phase)}, targeting \texttt{"requirements elicitation"}, \texttt{"requirements validation"}, \texttt{"requirements discovery"}, and \texttt{"requirements engineering"}; and (3) \textit{Facet 3 (Phenomenon / Defect)}, encompassing \texttt{"missing requirements"}, \texttt{"unclear requirements"}, \texttt{"ambiguity"}, \texttt{"completeness"}, \texttt{"tacit knowledge"}, and \texttt{"specification defect"}.

\subsection{Inclusion and exclusion criteria}
To guarantee methodological rigor and eliminate non-empirical or out-of-scope publications, studies were evaluated against the criteria detailed in Table~\ref{tab:inclusion_exclusion}.

\begin{table}[htbp]
\centering
\caption{Inclusion and Exclusion Criteria}
\label{tab:inclusion_exclusion}
\begin{tabularx}{\columnwidth}{l X}
\toprule
\textbf{Code} & \textbf{Criterion Description} \\
\midrule
\multicolumn{2}{l}{\textit{\textbf{Inclusion Criteria (IC)}}} \\
\textbf{IC1} & Peer-reviewed journal articles, conference proceedings, or workshop papers published in recognized SE and RE venues. \\
\textbf{IC2} & Empirical investigation (controlled experiments, quasi-experiments, multi-case studies, or mapping studies) examining prototyping in requirements elicitation or validation. \\
\textbf{IC3} & Direct evaluation of requirement quality attributes, specifically regarding missing requirements discovery, ambiguity reduction, or requirement typology. \\
\midrule
\multicolumn{2}{l}{\textit{\textbf{Exclusion Criteria (EC)}}} \\
\textbf{EC1} & Publications where prototyping refers solely to rapid software coding, architectural prototyping, or hardware emulation without requirements engineering stakeholder involvement. \\
\textbf{EC2} & Short vision papers ($< 4$ pages), extended abstracts, keynotes, non-peer-reviewed industry white papers, or tertiary overviews. \\
\textbf{EC3} & Studies lacking transparent empirical data collection, reproducible methodology, or clear defect reporting. \\
\bottomrule
\end{tabularx}
\end{table}

\subsection{Study selection protocol and selection funnel}
The selection process was conducted through a multi-stage screening funnel to ensure reproducibility. Table~\ref{tab:selection_funnel} outlines the numerical progression across the screening phases.

\begin{table}[htbp]
\centering
\caption{Study Selection Funnel Across Electronic Databases}
\label{tab:selection_funnel}
\begin{tabularx}{\columnwidth}{l X r}
\toprule
\textbf{Stage} & \textbf{Screening Activity} & \textbf{Yield} \\
\midrule
\textbf{Stage 1} & Initial keyword execution across 4 databases & 312 \\
                 & \quad $\bullet$ IEEE Xplore: 98 & \\
                 & \quad $\bullet$ ACM Digital Library: 74 & \\
                 & \quad $\bullet$ ScienceDirect: 68 & \\
                 & \quad $\bullet$ SpringerLink: 72 & \\
\textbf{Stage 2} & Automated deduplication & 224 \\
\textbf{Stage 3} & Title and abstract screening against IC/EC & 54 \\
\textbf{Stage 4} & Full-text detailed quality appraisal & 26 \\
\textbf{Stage 5} & Final Included Corpus (Core Empirical + Foundational) & \textbf{18} \\
\bottomrule
\end{tabularx}
\end{table}

The resulting corpus of 18 studies is partitioned into two mutually reinforcing tiers: (1) \textit{10 Core Primary Studies (2020--2025)}, comprising top-tier empirical investigations providing direct quantitative and qualitative measurements addressing RQ6.1, RQ6.2, and RQ6.3 \cite{bjarnason2023empirically, rueda2020requirements, uv2023empirical, rej2022evolve, antonelli2022iterative, bano2023videos, gupta2022strategic, acm2022proposal, diaz2021gender, cascon2025ambiguity}; and (2) \textit{8 Foundational and Seminal Studies}, establishing groundwork theoretical baselines for IREB standards, participative prototyping, elicitation maturity, interview ambiguity, and video scenarios \cite{pohl2016requirements, hickey1998prototyping, abad2018loud, jserd2019protopersona, davis2018maturity, ferrari2016ambiguity, spoletini2018interview, brill2010videos}.

\subsection{Quality appraisal and data extraction protocol}
Each eligible full-text article was assessed across four Quality Appraisal (QA) benchmarks: (QA1) Explicit formulation of research objectives; (QA2) Rigor and transparency of experimental/case design; (QA3) Objective measurement of requirement defect discovery or clarity; and (QA4) Validity threat discussion.

Data extraction was executed using a standardized protocol structured across five operational dimensions: (1) \textit{Bibliographic Profile}, recording authors, publication year, venue, and empirical design; (2) \textit{PAM Model Classification}, categorizing studies along the five dimensions (Purpose, Scope, Media, Exploration Strategy, Use) \cite{bjarnason2023empirically}; (3) \textit{Empirical Findings on Missing/Unclear Requirements}, extracting statistical indicators ($p$-values, effect sizes $\eta^2$), requirement modification percentages, and observed defect types; (4) \textit{Requirement Typology Susceptibility}, analyzing Functional (FR), Non-Functional (NFR), and Usability/GUI requirements; and (5) \textit{Tripartite Analytical Mapping}, synthesizing evidence across Rationale (\textit{Why}: cognitive and linguistic barriers), Operational Protocol (\textit{How}: interaction workflows), and Taxonomical Dimensions (\textit{Classification}: fidelity and scope).
